\documentclass{article}

\usepackage{amsmath,amssymb}
\usepackage{xurl}
\usepackage[colorlinks=true,linkcolor=blue,citecolor=blue,urlcolor=blue]{hyperref}
\setlength{\emergencystretch}{2em}

% Shared commands for notes in docs/paper-gaps/.

\DeclareUnicodeCharacter{2080}{\ifmmode{}_{0}\else\textsubscript{0}\fi}
\DeclareUnicodeCharacter{2081}{\ifmmode{}_{1}\else\textsubscript{1}\fi}
\DeclareUnicodeCharacter{2082}{\ifmmode{}_{2}\else\textsubscript{2}\fi}
\DeclareUnicodeCharacter{2099}{\ifmmode{}_{n}\else\textsubscript{n}\fi}

\newcommand{\Lean}{\textsc{Lean}}
\ExplSyntaxOn
\NewDocumentCommand{\leanid}{m}{
  \ifmmode
    \textsf{\detokenize{#1}}
  \else
    \group_begin:
    \urlstyle{sf}
    \tl_set:Nn \l_tmpa_tl {#1}
    \tl_replace_all:Nnn \l_tmpa_tl {\_} {_}
    \exp_args:NV \path \l_tmpa_tl
    \group_end:
  \fi
}
\ExplSyntaxOff
\providecommand{\tr}{\operatorname{tr}}
\newcommand{\tnleanIssueBase}{https://github.com/LionSR/TNLean/issues/}
\newcommand{\tnleanPrBase}{https://github.com/LionSR/TNLean/pull/}
\newcommand{\tnleanDocBase}{https://sirui-lu.com/TNLean/docs/}
\newcommand{\ghissue}[1]{\href{\tnleanIssueBase#1}{GitHub issue \##1}}
\newcommand{\ghpr}[1]{\href{\tnleanPrBase#1}{PR \##1}}
\newcommand{\doclink}[1]{\href{\tnleanDocBase#1.html}{\path{#1}}}

% Dirac notation and the identity operator, as in blueprint/src/macros/common.tex.
% \providecommand keeps note-local definitions authoritative.
\usepackage{dsfont}
\providecommand{\ket}[1]{|#1\rangle}
\providecommand{\bra}[1]{\langle #1|}
\providecommand{\braket}[2]{\langle #1 | #2 \rangle}
\providecommand{\ketbra}[2]{|#1\rangle\!\langle #2|}
\providecommand{\Id}{\mathds{1}}
\providecommand{\one}{\mathds{1}}
\providecommand{\C}{\mathbb{C}}
\providecommand{\MN}[1]{M_{#1}(\C)}
\providecommand{\MD}{M_D(\C)}

% External referencing for paper sources.  The build helper writes sanitized
% auxiliary files under build/paper-aux/ and excludes duplicated source labels.
\usepackage{xr}

% Import a generated paper auxiliary file with a caller-supplied label prefix.
% The candidate paths support builds from docs/paper-gaps/, the repository root,
% and build/paper-aux/ itself.
\newcommand{\importpaperaux}[2]{%
  \IfFileExists{../../build/paper-aux/#2.aux}{%
    \externaldocument[#1]{../../build/paper-aux/#2}%
  }{%
    \IfFileExists{build/paper-aux/#2.aux}{%
      \externaldocument[#1]{build/paper-aux/#2}%
    }{%
      \IfFileExists{#2.aux}{%
        \externaldocument[#1]{#2}%
      }{}%
    }%
  }%
}

% General external reference macros
\newcommand{\paperref}[2]{\ref{#1-#2}}
\newcommand{\papereqref}[2]{(\ref{#1-#2})}

% CPSV16 source (1606.00608 / MPDO-22-12-17-2.tex) external document and shortcuts
\importpaperaux{cpsv16-}{1606.00608/MPDO-22-12-17-2-tnlean}
\newcommand{\cpsvref}[1]{\paperref{cpsv16}{#1}}
\newcommand{\cpsveqref}[1]{\papereqref{cpsv16}{#1}}

% Verdict marker: \gapnote{<kind>}{<status>}. Kind is the policy
% classification (clarification, local-correction, scope-restriction,
% unfaithful, false-source, open-gap); status is open, wip (elimination
% underway), resolved, or historical. Machine-read by the site index;
% prints a verdict line.
\newcommand{\gapnote}[2]{\par\noindent\textbf{Verdict:} #1 (#2).\par}


% Fallbacks keep this author document compilable when it is built outside its
% source directory and a different command.tex is found first.
\providecommand{\Lean}{\textsc{Lean}}
\providecommand{\leanid}[1]{\textsf{\detokenize{#1}}}
\providecommand{\doclink}[1]{%
    \href{https://sirui-lu.com/TNLean/docs/#1.html}{\path{#1}}}
\providecommand{\ghpr}[1]{%
    \href{https://github.com/LionSR/TNLean/pull/#1}{PR \##1}}
\providecommand{\gapnote}[2]{%
    \par\noindent\textbf{Verdict:} #1 (#2).\par}
\providecommand{\tr}{\operatorname{tr}}
\providecommand{\ket}[1]{|#1\rangle}

\title{Per-block RFP Isometry Form and the Source Isometry Condition in CPSV16}
\author{}
\date{2026-07-31}

\begin{document}
\maketitle
\gapnote{open-gap}{open}

\section{Source statement}

The pure-state RFP characterization in Section~\cpsvref{Sec:MPS} of
Ref.~\cite{Cirac2016MPDO_arXiv} writes a canonical-form RFP tensor as
\begin{align}
    A^i
        &=
    \bigoplus_{j=1}^g\bigoplus_{q=1}^{r_j}
    \mu_{j,q}X_{j,q}\Lambda_jU_j^iX_{j,q}^{-1}.
    \label{eq:cpsv16_canonical_rfp_tensor}
\end{align}
This is Equation~(\cpsvref{III_CFI_RFP}) of Ref.~\cite{Cirac2016MPDO_arXiv},
where the matrices $\Lambda_j$ are diagonal, positive, and trace-normalized. The
same theorem imposes the joint isometry condition
\begin{align}
    \sum_{i=1}^s
    (U_j^i)_{\alpha,\beta}
    \overline{(U_{j'}^i)_{\alpha',\beta'}}
        &=
    \delta_{j,j'}\delta_{\alpha,\alpha'}\delta_{\beta,\beta'},
    \label{eq:cpsv16_source_joint_isometry}
\end{align}
which is Equation~(\cpsvref{eq:III_isometry}) of
Ref.~\cite{Cirac2016MPDO_arXiv}. The diagonal case $j=j'$ asserts full
pair-index orthonormality in the virtual indices $(\alpha,\beta)$, while the
Kronecker delta $\delta_{j,j'}$ adds cross-block orthogonality: for $j\ne j'$,
the physical-index isometries associated with distinct normal-tensor blocks are
mutually orthogonal. Corollary~\cpsvref{III_cor3} of
Ref.~\cite{Cirac2016MPDO_arXiv} invokes the same isometry condition for the
BNT elements.\footnote{Local source:
\path{Papers/1606.00608/MPDO-22-12-17-2.tex:543--554} and
\path{Papers/1606.00608/MPDO-22-12-17-2.tex:583--589}.}

\section{Formal boundary}

The single-normal-tensor characterization is formalized under the source
hypotheses. When $A$ is normal in the sense of Ref.~\cite{Cirac2016MPDO_arXiv},
its bond dimension is positive,\footnote{Formal declaration:
\leanid{MPSTensor.IsNormalTensor.bondDim_ne_zero}.} and
\begin{align}
    A\text{ is a renormalization fixed point}
        \quad\Longleftrightarrow\quad
    A\text{ is in the square-root isometry canonical form}.
    \label{eq:cpsv16_normal_rfp_iff_sqrt_isometry}
\end{align}
The forward implication in Eq.~\ref{eq:cpsv16_normal_rfp_iff_sqrt_isometry}
passes to the trace-preserving Perron gauge, applies the left-canonical
structural theorem there, and pulls the result back along the similarity; the
converse follows directly from the rank-one transfer formula.\footnote{Formal
declarations:
\leanid{MPSTensor.GaugeEquiv.isIsometryCanonicalForm_iff} in
\doclink{TNLean/MPS/RFP/StructuralFull} and
\leanid{MPSTensor.IsNormalTensor.isTransferIdempotent_iff_isIsometryCanonicalForm}
in \doclink{TNLean/MPS/RFP/NormalIsometryCharacterization}.}

Applied to each block of an indexed family, the same theorem gives the
following.\footnote{Formal declaration:
\leanid{MPSTensor.rfp_nt_structural_full} in
\doclink{TNLean/MPS/RFP/StructuralFull}.}
For each block $k$, let $D_k$ denote its bond dimension. The theorem gives
\begin{align}
    A_k^i
        &=
    X_k\Lambda_kU_k^iX_k^{-1},
    \label{eq:cpsv16_per_block_decomposition}\\
    \sum_i (U_k^i)^\dagger U_k^i
        &=
    I,
    \label{eq:cpsv16_per_block_contracted_identity}\\
    \sum_i
    \overline{(U_k^i)_{\alpha,\beta}}
    (U_k^i)_{\alpha',\beta'}
        &=
    D_k^{-1}\delta_{\alpha,\alpha'}\delta_{\beta,\beta'}.
    \label{eq:cpsv16_per_block_normalized_pair_isometry}
\end{align}
Here and below, $I$ denotes the identity matrix of the appropriate dimension.
Equation~\ref{eq:cpsv16_per_block_contracted_identity} is the contracted
identity
\begin{align}
    \sum_{i,\alpha}
    \overline{(U_k^i)_{\alpha,\beta}}
    (U_k^i)_{\alpha,\beta'}
        &=
    \delta_{\beta,\beta'}.
    \label{eq:cpsv16_per_block_contracted_entries}
\end{align}
Equation~\ref{eq:cpsv16_per_block_normalized_pair_isometry} records diagonal
pair-index orthonormality in the normalization of the formal theorem. Compared
with Eq.~\ref{eq:cpsv16_source_joint_isometry}, the entries of $U_k^i$ carry the
normalized matrix-unit factor $D_k^{-1/2}$, making
Eq.~\ref{eq:cpsv16_per_block_normalized_pair_isometry} the matching pair-index
relation. The normalization comparison with
Equation~(\cpsvref{eq:III_isometry}) of Ref.~\cite{Cirac2016MPDO_arXiv} is thus
explicit.

A comparison theorem rescales this single-block decomposition by
\begin{align}
    U_k^i
        &\longmapsto
    D_k^{1/2}U_k^i,
    &
    \Lambda_k
        &\longmapsto
    D_k^{-1/2}\Lambda_k.
    \label{eq:cpsv16_unit_pair_rescaling}
\end{align}
\footnote{Formal declaration:
\leanid{MPSTensor.rfp_nt_structural_full_unit_pair} in
\doclink{TNLean/MPS/RFP/StructuralFull}.}
This rescaling leaves $\Lambda_kU_k^i$ unchanged and recovers the diagonal
pair-index equation in the source unit convention,
\begin{align}
    \sum_i
    \overline{(U_k^i)_{\alpha,\beta}}
    (U_k^i)_{\alpha',\beta'}
        &=
    \delta_{\alpha,\alpha'}\delta_{\beta,\beta'}.
    \label{eq:cpsv16_per_block_unit_pair_isometry}
\end{align}
Both conventions are available, and their relation is fully recorded.

The structural datum in Appendix~\cpsvref{AppendixPure} used in the conditional
parent-Hamiltonian theorem adopts the unit pair-index convention, so the
coefficient problem starts from the source normalization. The same appendix of
Ref.~\cite{Cirac2016MPDO_arXiv} gives the coefficient form of the basic-vector
expression as
\begin{align}
    C_L(\sigma)
        &=
    \sum_{\alpha_0,\ldots,\alpha_{L-1}}
    \prod_{t=0}^{L-1}
    \Lambda_{\alpha_t}
    (U^{\sigma_t})_{\alpha_t,\alpha_{t+1}},
    \qquad
    \alpha_L=\alpha_0.
    \label{eq:cpsv16_cyclic_coefficient}
\end{align}
This is the coefficient form of
\begin{align}
    \ket{V^{(L)}(A_j)}
        &=
    U^{\otimes L}\ket{\varphi_j}^{\otimes L},
    \label{eq:cpsv16_basic_vector_expression}
\end{align}
where $\varphi_j$ is shared between $b_t$ and $a_{t+1}$. The even-chain
physical-pair factorization considered in the conditional parent-Hamiltonian
theorem is therefore a separate statement rather than the source expression
itself.

For every nonempty chain length and chosen structural datum, the formal
development proves the cyclic coefficient formula for the original tensor
after the basis change of Appendix~\cpsvref{AppendixPure}:
\begin{align}
    V^{(L)}(A)_\sigma
        &=
    C_L(\sigma).
    \label{eq:cpsv16_formal_cyclic_virtual_pair_state}
\end{align}
\footnote{Formal theorem:
\leanid{MPSTensor.AppendixBStructuralData.mpv_eq_cyclicVirtualPairState}.} This
cyclic identity does not by itself establish the even-chain physical-pair
factorization.

The pair-index isometry in Eq.~\ref{eq:cpsv16_per_block_unit_pair_isometry}
sums over the physical letter $i$ and thus does not supply a coefficient identity
for an individual physical word $\sigma$. Proving a fixed-word physical-pair
factorization requires the source basic-vector expression in
Eq.~\ref{eq:cpsv16_basic_vector_expression} or an explicit projector
construction, because the physical indices in $C_L(\sigma)$ are fixed rather
than summed.

This obstruction is structural rather than a formal artifact. Consider two
virtual indices, physical labels $i=(a,b)$, and
\begin{align}
    (U^{(a,b)})_{x,y}
        &=
    \delta_{x,a}\delta_{y,b},
    &
    A^{(a,b)}
        &=
    \lambda_aE_{a,b},
    \qquad
    \lambda_0,\lambda_1>0.
    \label{eq:cpsv16_fixed_word_counterexample_tensor}
\end{align}
The unit pair-index isometry holds. For the length-four word
\begin{align}
    \sigma
        &=
    ((0,1),(1,0),(1,0),(0,1)),
    \label{eq:cpsv16_fixed_word_counterexample_word}
\end{align}
the cyclic coefficient vanishes because virtual indices fail to match:
\begin{align}
    C_4(\sigma)
        &=
    0.
    \label{eq:cpsv16_counterexample_cyclic_coefficient}
\end{align}
The physical-pair product, however, evaluates to
\begin{align}
    (\lambda_0\lambda_1)(\lambda_1\lambda_0)
        &=
    \lambda_0^2\lambda_1^2
        \ne
    0.
    \label{eq:cpsv16_counterexample_physical_pair_product}
\end{align}
Equations~\ref{eq:cpsv16_counterexample_cyclic_coefficient}
and~\ref{eq:cpsv16_counterexample_physical_pair_product} show that the
even-chain physical-pair coefficient identity cannot be deduced from the unit
pair-index isometry alone.

The per-block isometry form does not constrain $U_k$ and $U_\ell$ for
$k\ne\ell$. The cross-block relations
\begin{align}
    \sum_i
    (U_k^i)_{\alpha,\beta}
    \overline{(U_\ell^i)_{\alpha',\beta'}}
        &=
    0,
    \qquad
    k\ne\ell,
    \label{eq:cpsv16_cross_block_residual_isometry}
\end{align}
form the $j\ne j'$ portion of Eq.~\ref{eq:cpsv16_source_joint_isometry}; they
are proved at the normal-tensor-block level in
Section~\ref{sec:cpsv16_cross_block}.

\section{Square-root diagonal (local correction)}
\label{sec:cpsv16_square_root_correction}

Equation~(\cpsvref{eq_III_NT_RFP}) in Ref.~\cite{Cirac2016MPDO_arXiv} displays
the normal-tensor decomposition as
\begin{align}
    A^i
        &=
    X\Lambda U^iX^{-1},
    \label{eq:cpsv16_source_normal_tensor_form}
\end{align}
combining the unsquared diagonal matrix $\Lambda$ with the unit isometry
condition\footnote{Local source:
\path{Papers/1606.00608/MPDO-22-12-17-2.tex:1278} for the displayed equation and
\path{Papers/1606.00608/MPDO-22-12-17-2.tex:1281--1283} for the unit isometry.}
\begin{align}
    \sum_i
    (U^i)_{\alpha,\beta}
    \overline{(U^i)_{\alpha',\beta'}}
        &=
    \delta_{\alpha,\alpha'}\delta_{\beta,\beta'}.
    \label{eq:cpsv16_source_normal_tensor_unit_isometry}
\end{align}
Taken literally, these two conditions are incompatible. The proof in
Ref.~\cite{Cirac2016MPDO_arXiv} builds the reference tensor\footnote{Local
source: \path{Papers/1606.00608/MPDO-22-12-17-2.tex:1300}.}
\begin{align}
    \widehat A^{(\alpha',\beta')}_{\alpha,\beta}
        &=
    \delta_{\alpha,\alpha'}\delta_{\beta,\beta'}
    \sqrt{\Lambda_\alpha},
    \label{eq:cpsv16_source_reference_tensor}
\end{align}
which is left-canonical with fixed point $\operatorname{diag}(\Lambda)$ and
$\tr(\Lambda)=1$. Setting $\widehat A=\Lambda U$ forces
\begin{align}
    U^{(\alpha',\beta')}_{\alpha,\beta}
        &=
    \frac{\delta_{\alpha,\alpha'}\delta_{\beta,\beta'}}
         {\sqrt{\Lambda_\alpha}},
    \label{eq:cpsv16_residual_after_full_diagonal}
\end{align}
giving the pair-index sum
\begin{align}
    \sum_i
    (U^i)_{\alpha,\beta}
    \overline{(U^i)_{\alpha',\beta'}}
        &=
    \frac{\delta_{\alpha,\alpha'}\delta_{\beta,\beta'}}
         {\Lambda_\alpha}
        \ne
    \delta_{\alpha,\alpha'}\delta_{\beta,\beta'},
    \qquad
    \Lambda_\alpha\ne1.
    \label{eq:cpsv16_full_diagonal_nonunit_residual}
\end{align}
Whenever $\Lambda_\alpha\ne1$, this residual is not a unit isometry. Setting
$\widehat A=\sqrt{\Lambda}\,U$ yields instead
\begin{align}
    U^{(\alpha',\beta')}_{\alpha,\beta}
        &=
    \delta_{\alpha,\alpha'}\delta_{\beta,\beta'},
    \label{eq:cpsv16_residual_after_square_root_diagonal}
\end{align}
with pair-index sum
\begin{align}
    \sum_i
    (U^i)_{\alpha,\beta}
    \overline{(U^i)_{\alpha',\beta'}}
        &=
    \delta_{\alpha,\alpha'}\delta_{\beta,\beta'},
    \label{eq:cpsv16_square_root_unit_residual}
\end{align}
which is indeed a unit isometry.

The formalization therefore adopts the square-root diagonal decomposition
\begin{align}
    A^i
        &=
    X\sqrt{\Lambda}\,U^iX^{-1},
    \label{eq:cpsv16_corrected_square_root_form}
\end{align}
preserving both the trace normalization $\tr(\Lambda)=1$ on $\Lambda$ and the
unit isometry condition on $U$. The replacement $\Lambda\mapsto\sqrt{\Lambda}$
is a local correction to the printed display. The source itself uses
$\sqrt{\Lambda_\alpha}$ in the reference tensor at line~1300, confirming that
$\Lambda$ at line~1278 is an imprecise notation for the dressing.

\section{Trace normalization (discharged)}
\label{sec:cpsv16_trace_normalization}

The trace-normalized isometry canonical form is captured by a dedicated
predicate.\footnote{Formal declaration: \leanid{MPSTensor.IsIsometryCanonicalForm}
in \doclink{TNLean/MPS/RFP/StructuralFull}.}
It asserts the existence of an invertible matrix $X$, a positive diagonal
matrix $\Lambda$ with $\sum_\alpha\Lambda_\alpha=1$, and a unit pair-index
isometry $U$ such that
\begin{align}
    A^i
        &=
    X\sqrt{\Lambda}\,U^iX^{-1},
    \label{eq:cpsv16_trace_normalized_sqrt_decomposition}\\
    \sum_i
    \overline{(U^i)_{\alpha,\beta}}
    (U^i)_{\alpha',\beta'}
        &=
    \delta_{\alpha,\alpha'}\delta_{\beta,\beta'}.
    \label{eq:cpsv16_trace_normalized_unit_isometry}
\end{align}
The square root on the diagonal matches the reference tensor in
Section~\ref{sec:cpsv16_square_root_correction}. The trace normalization
applies directly to $\Lambda$, or equivalently to the normalized diagonal fixed
point $\rho=\operatorname{diag}(\Lambda)$:
\begin{align}
    \Lambda_\alpha
        &=
    \frac{\rho_{\alpha,\alpha}}{\tr\rho},
    &
    \sum_\alpha\Lambda_\alpha
        &=
    1,
    \label{eq:cpsv16_fixed_point_trace_normalization}
\end{align}
where the normalization sum follows from
\begin{align}
    \sum_\alpha\rho_{\alpha,\alpha}
        &=
    \tr\rho.
    \label{eq:cpsv16_diagonal_fixed_point_trace}
\end{align}

The forward theorem establishes this form for the Perron representative in
left-canonical gauge.\footnote{Formal theorem:
\leanid{MPSTensor.isIsometryCanonicalForm_of_rfp_nt}.}
By gauge invariance, every normal tensor is a renormalization fixed point if
and only if it admits the square-root form in
Eq.~\ref{eq:cpsv16_trace_normalized_sqrt_decomposition}, without requiring
left-canonical normalization or bond-dimension restrictions on the original
tensor.\footnote{Formal equivalence:
\leanid{MPSTensor.IsNormalTensor.isTransferIdempotent_iff_isIsometryCanonicalForm}.}

This discharges the single-normal-tensor characterization and its trace
normalization. The cross-block orthogonality ($\delta_{j,j'}$) between
distinct normal-tensor blocks is treated next.

\section{Cross-block residual isometry}
\label{sec:cpsv16_cross_block}

Cross-block orthogonality ($\delta_{j,j'}$) is formalized at the level of
normal-tensor blocks. Let $(A_j)_j$ be a family of normal, irreducible,
left-canonical blocks of positive dimensions, pairwise inequivalent under gauge
transformations and phase shifts, such that $\bigoplus_jA_j$ is a
renormalization fixed point. For distinct blocks, the mixed transfer operator
\begin{align}
    \mathcal E_{j,j'}
        &=
    \sum_i A_j^i\otimes\overline{A_{j'}^i}
    \label{eq:cpsv16_mixed_transfer_definition}
\end{align}
vanishes identically:
\begin{align}
    \mathcal E_{j,j'}
        &=
    0,
    \qquad
    j\ne j'.
    \label{eq:cpsv16_mixed_transfer_vanishing}
\end{align}
\footnote{Formal theorem:
\leanid{MPSTensor.isBNTLocallyOrthogonal_of_isTransferIdempotent_directSum}.}
The proof decomposes the direct-sum transfer map block by block.\footnote{Formal
decomposition: \leanid{MPSTensor.blockDiagonal'_transferSum_toBlock}.}
Transfer-map idempotence of the direct sum implies that each $\mathcal
E_{j,j'}$ is idempotent; because distinct blocks have peripheral spectral
radius strictly below one, every such idempotent must vanish.

Passing to the residual isometries uses the canonical form of each block.
Writing
\begin{align}
    A_j^i
        &=
    X_j\sqrt{\Lambda_j}\,U_j^iX_j^{-1},
    \label{eq:cpsv16_cross_block_sqrt_decomposition}
\end{align}
with invertible $X_j$ and $\sqrt{\Lambda_j}$, the mixed transfer operator
transforms covariantly under two-sided gauge changes:
\begin{align}
    \mathcal E_{A_j,A_{j'}}(Y)
        &=
    X_j\sqrt{\Lambda_j}
    \mathcal E_{U_j,U_{j'}}
    (X_j^{-1}Y(X_{j'}^{-1})^\dagger)
    (X_{j'}\sqrt{\Lambda_{j'}})^\dagger.
    \label{eq:cpsv16_mixed_transfer_covariance}
\end{align}
\footnote{Formal covariance: \leanid{MPSTensor.mixedMapLM_conj_apply}.}
Cancelling the invertible outer factors yields
\begin{align}
    \mathcal E_{A_j,A_{j'}}
        =
    0
    \quad\Longrightarrow\quad
    \mathcal E_{U_j,U_{j'}}
        &=
    0.
    \label{eq:cpsv16_cancel_mixed_transfer_gauges}
\end{align}
Evaluating Eq.~\ref{eq:cpsv16_cancel_mixed_transfer_gauges} on matrix units
gives
\begin{align}
    \sum_i
    (U_j^i)_{\alpha,\beta}
    \overline{(U_{j'}^i)_{\alpha',\beta'}}
        &=
    0,
    \qquad
    j\ne j'.
    \label{eq:cpsv16_residual_cross_block_entries}
\end{align}
Combined with within-block orthonormality,
Eq.~\ref{eq:cpsv16_residual_cross_block_entries} establishes the full source
condition in Eq.~\ref{eq:cpsv16_source_joint_isometry}. The formal development
packages this condition as a residual isometry family\footnote{Formal
predicate: \leanid{MPSTensor.IsResidualIsometryFamily}.} and proves it for
direct sums with idempotent transfer map,\footnote{Formal theorem:
\leanid{MPSTensor.exists_residualIsometryFamily_of_isTransferIdempotent_directSum}.}
discharging the residual-isometry claim of Corollary~\cpsvref{III_cor3} in
Ref.~\cite{Cirac2016MPDO_arXiv} at the block level.

\section{Corollary 3.12: phase-class representatives}
\label{sec:cpsv16_phase_class_representatives}

Corollary~\cpsvref{III_cor3} of Ref.~\cite{Cirac2016MPDO_arXiv} applies this
joint residual-isometry conclusion to a basis of normal tensors for a
canonical-form renormalization fixed point. Under the nonzero-coefficient
convention of
\path{docs/paper-gaps/cpsv16_bnt_uniqueness_zero_coefficient.tex}, each basis
element is gauge-phase equivalent to a canonical-form block, reducing the
statement to phase-class representatives.

Write
\begin{align}
    A^i
        &=
    U^\dagger\left(\bigoplus_k\mu_kA_k^i\right)U,
    &
    UU^\dagger
        &=
    I,
    \label{eq:cpsv16_canonical_coisometric_decomposition}
\end{align}
where each $A_k$ is normal and each $\mu_k$ is nonzero. Choosing one
representative $C_j$ per phase class, if $A$ has an idempotent transfer map,
there exist invertible matrices $X_j$, positive weights $\lambda_{j,\alpha}$
with $\sum_\alpha\lambda_{j,\alpha}=1$, and tensors $V_j$ such that
\begin{align}
    C_j^i
        &=
    X_j\operatorname{diag}
    \left(\sqrt{\lambda_{j,\alpha}}\right)V_j^iX_j^{-1},
    \label{eq:cpsv16_phase_representative_sqrt_form}\\
    \sum_i
    (V_j^i)_{\alpha,\beta}
    \overline{(V_{j'}^i)_{\alpha',\beta'}}
        &=
    \delta_{j,j'}\delta_{\alpha,\alpha'}\delta_{\beta,\beta'}.
    \label{eq:cpsv16_phase_representative_joint_isometry}
\end{align}
This holds also when the index set of phase classes is empty. The result is
formalized for phase-class representatives.\footnote{Formal declarations:
\leanid{MPSTensor.CPSVCanonicalFormData.BNTRefinement.exists_residualIsometryFamily_of_isTransferIdempotent}
and predicate wrapper
\leanid{MPSTensor.IsCPSVCanonicalForm.exists_bntRefinement_residualIsometryFamily_of_isTransferIdempotent}
in \doclink{TNLean/MPS/RFP/CPSVCanonicalForm}.}

The proof recovers the weighted direct sum through the ambient coisometry.
For each block, the weight has modulus one and the block is an RFP:\footnote{Formal
declaration:
\leanid{MPSTensor.CPSVCanonicalFormData.weight_norm_one_and_block_rfp}.}
\begin{align}
    \lvert\mu_j\rvert
        &=
    1,
    &
    \mathcal E_{C_j}^2
        &=
    \mathcal E_{C_j},
    \label{eq:cpsv16_weight_norm_and_block_rfp}
\end{align}
and the single-block theorem provides the square-root decomposition. Putting
each representative into a trace-preserving Perron gauge and tracking its
unit-modulus weight, idempotence of the direct sum restricts to each mixed pair.
Because distinct phase-class representatives are not gauge-phase equivalent,
the mixed spectral gap forces
\begin{align}
    \mathcal E_{\mu_jC_j,\mu_{j'}C_{j'}}
        &=
    0,
    \qquad
    j\ne j'.
    \label{eq:cpsv16_weighted_phase_mixed_transfer_zero}
\end{align}
Cancelling the nonzero scalar factors and invertible gauges yields
\begin{align}
    \mathcal E_{C_j,C_{j'}}
        &=
    0,
    \qquad
    j\ne j'.
    \label{eq:cpsv16_phase_mixed_transfer_zero}
\end{align}
Factoring out the positive diagonal matrices and evaluating on matrix units
gives Eq.~\ref{eq:cpsv16_phase_representative_joint_isometry}. The off-diagonal
equations require no extra direct-sum assumptions, establishing
Corollary~\cpsvref{III_cor3} up to the square-root correction of
Section~\ref{sec:cpsv16_square_root_correction}.

\section{Backward direction (multiplicity-one)}
\label{sec:cpsv16_backward_multiplicity_one}

The converse implication is likewise formalized at the block level. If each
block $B_k$ is in isometry canonical form and all off-diagonal mixed transfer
operators vanish,
\begin{align}
    \mathcal E_{j,j'}
        &=
    0,
    \qquad
    j\ne j',
    \label{eq:cpsv16_backward_cross_block_hypothesis}
\end{align}
then the direct sum is a renormalization fixed point:
\begin{align}
    \bigoplus_kB_k
        &\text{ is a renormalization fixed point}.
    \label{eq:cpsv16_backward_direct_sum_rfp}
\end{align}
\footnote{Formal declaration:
\leanid{MPSTensor.isTransferIdempotent_directSumTensor_of_isIsometryCanonicalForm}
in \doclink{TNLean/MPS/RFP/ResidualIsometry}.}
This follows from transfer-map block additivity.\footnote{Formal theorem:
\leanid{MPSTensor.blockTransferSum_idempotent_of_isIsometryCanonicalForm}.}
Decomposing the direct-sum transfer map blockwise,\footnote{Formal decomposition:
\leanid{MPSTensor.blockDiagonal'_transferSum_toBlock}.} the diagonal blocks are
idempotent,\footnote{Formal theorem:
\leanid{MPSTensor.isTransferIdempotent_of_isIsometryCanonicalForm}.} and the
off-diagonal blocks vanish by Eq.~\ref{eq:cpsv16_backward_cross_block_hypothesis}.
Because this off-diagonal condition is the $j\ne j'$ clause of
Eq.~\ref{eq:cpsv16_source_joint_isometry}, it is part of the source isometry
form rather than an extraneous hypothesis.

This covers the distinct-block, multiplicity-one, unit-phase case of
Eq.~\ref{eq:cpsv16_canonical_rfp_tensor}, where the inner sum over copies
$q=1,\ldots,r_j$ is trivial and every phase $\mu_{j,q}$ equals $1$. The
full-multiplicity converse includes repeated copies of each block with distinct
phases and scalar weights. Treating those copies requires care with the physical
direct sum indexed by $q$, since a naive virtual-block substitution fails the
defining physical-isometry condition.

For a sector decomposition satisfying the BNT canonical-form predicate,
explicit blockwise hypotheses are no longer required. The basis-direct-sum
equivalence extracts positive bond dimensions, normality, irreducibility,
left-canonical normalization, and gauge-phase distinctness directly from the
predicate.\footnote{Formal theorem:
\leanid{MPSTensor.IsBNTCanonicalForm.isTransferIdempotent_basisDirectSum_iff}.}
In the forward direction, the trace-normalized decompositions and the residual
isometry family are supplied unconditionally.\footnote{Formal theorem:
\leanid{MPSTensor.IsBNTCanonicalForm.exists_residualIsometryFamily_of_isTransferIdempotent_basisDirectSum}.}
This settles the canonical-form extraction for the multiplicity-one basis direct
sum, though it does not identify that sum with the repeated-copy tensor carrying
weights $\mu_{j,q}$.

For an arbitrary block direct sum $\bigoplus_kB_k$, whole-tensor transfer
idempotence reduces to pairwise block equations:
\begin{align}
    \mathcal E_{\bigoplus_kB_k}^2
        =
    \mathcal E_{\bigoplus_kB_k}
    \quad\Longleftrightarrow\quad
    \mathcal E_{j,j'}^2
        &=
    \mathcal E_{j,j'}
    \quad\text{for every }j,j'.
    \label{eq:cpsv16_direct_sum_pairwise_idempotence}
\end{align}
\footnote{Formal equivalence:
\leanid{MPSTensor.isTransferIdempotent_directSumTensor_iff_pairwise_mixedMapLM_isIdempotentElem}.}
Equation~\ref{eq:cpsv16_direct_sum_pairwise_idempotence} shows that repeated
virtual copies impose relations between distinct copies of the same normal
block; testing only representative blocks misses part of the whole-tensor
fixed-point condition.

\section{Why the fixed-letter virtual-block reading is excluded}
\label{sec:cpsv16_fixed_letter_exclusion}

Equation~\ref{eq:cpsv16_canonical_rfp_tensor} allows a normal tensor $j$ to
appear in multiple copies $q$ with independent unimodular phases $\mu_{j,q}$.
Interpreted as a standard virtual block diagonal for each fixed physical
letter, a single physical letter with scalar representative $B^0=(1)$ and copy
phases $1$ and $-1$ yields
\begin{align}
    A^0
        &=
    \operatorname{diag}(1,-1),
    \label{eq:cpsv16_literal_virtual_block_example}
\end{align}
which satisfies neither Eq.~(\cpsvref{AA=A}) nor transfer-map idempotence.

This reading, however, does not match the source. The text following
Equation~(\cpsvref{III_CFI_RFP}) in Ref.~\cite{Cirac2016MPDO_arXiv} explains
that $j$ and $q$ connect the physical and virtual indices, producing a direct
sum in both spaces up to the physical isometry.\footnote{Local source:
\path{Papers/1606.00608/MPDO-22-12-17-2.tex:559--563}.}
Because a single physical letter cannot host a nontrivial physical direct sum,
Eq.~\ref{eq:cpsv16_literal_virtual_block_example} falls outside the intended
class and does not constitute a counterexample to the source claim.

What the source omits is the coordinate data for routing $q$ physically. The
isometry display at lines~551--554 indexes $j$ without mentioning $q$, while
lines~561--563 describe physical and virtual direct sums without defining a
$q$-dependent isometry or stating its isometry relation. A rigorous construction
would also need dimension bounds to embed the direct sum into the physical space,
which the theorem does not provide. Introducing such hypotheses in \Lean{} would
alter the source statement. Formalizing the full repeated-copy theorem therefore
requires a corrected source formulation.

\section{Elimination plan}
\label{sec:cpsv16_elimination_plan}

Canonical-form extraction for the multiplicity-one basis direct sum is
complete, as established by the theorems in
Section~\ref{sec:cpsv16_backward_multiplicity_one}. For literal repeated
virtual block sums, the mixed transfer operator obeys the scaling law
\begin{align}
    \mathcal E_{\mu B,\nu C}
        &=
    \mu\overline{\nu} \mathcal E_{B,C},
    \label{eq:cpsv16_mixed_transfer_scaling}
\end{align}
\footnote{Formal identity: \leanid{Kraus.mixedMapLM_smul}.}
Combining this with the pairwise criterion in
Eq.~\ref{eq:cpsv16_direct_sum_pairwise_idempotence} shows that for repeated
copies of a nonzero idempotent block, whole-tensor idempotence forces
\begin{align}
    \mu_q\overline{\mu_{q'}}
        &=
    1,
    \label{eq:cpsv16_repeated_copy_phase_product}\\
    \mu_q
        &=
    \mu_{q'}
    \label{eq:cpsv16_repeated_copy_phase_equality}
\end{align}
and hence $\mu_q=\mu_{q'}$ for unimodular phases. This corrected literal-block
result is established unconditionally.\footnote{Formal theorem:
\leanid{MPSTensor.phases_eq_of_isTransferIdempotent_directSum_scaled_self}.} It
resolves the literal block construction by correcting the statement rather than
formalizing the printed independent-phase claim.

A formalization of the independent-phase claim requires an erratum specifying
the $q$-dependent physical map, its isometry relations, and the necessary
dimension constraints. These missing data cannot be supplied as ad hoc
hypotheses.

\section{Conclusion}

The per-block isometry form admits two normalization conventions, both
formalized explicitly. In the $D_k^{-1/2}$ convention, the residual tensor
satisfies
\begin{align}
    \sum_i(U_k^i)^\dagger U_k^i
        &=
    I,
    \label{eq:cpsv16_conclusion_normalized_contraction}
\end{align}
with pair-index right-hand side
$D_k^{-1}\delta_{\alpha,\alpha'}\delta_{\beta,\beta'}$. In the rescaled unit
pair-index convention, it satisfies
\begin{align}
    \sum_i
    \overline{(U_k^i)_{\alpha,\beta}}
    (U_k^i)_{\alpha',\beta'}
        &=
    \delta_{\alpha,\alpha'}\delta_{\beta,\beta'},
    \label{eq:cpsv16_conclusion_unit_pair_isometry}
\end{align}
while the contracted identity becomes
\begin{align}
    \sum_i(U_k^i)^\dagger U_k^i
        &=
    D_kI.
    \label{eq:cpsv16_conclusion_unit_pair_contraction}
\end{align}
The relation between these two scalings is therefore fully accounted for.

For every normal tensor, whose bond dimension is positive, transfer-map
idempotence is equivalent to the trace-normalized square-root form.\footnote{Formal
equivalence:
\leanid{MPSTensor.IsNormalTensor.isTransferIdempotent_iff_isIsometryCanonicalForm}.}
The diagonal weight is
\begin{align}
    \Lambda_\alpha
        &=
    \frac{\rho_{\alpha,\alpha}}{\tr\rho},
    \label{eq:cpsv16_conclusion_diagonal_weight}
\end{align}
and $\sqrt{\Lambda}$ dresses the unit isometry. This discharges the
single-normal-tensor lemma under its source hypotheses with the corrected
square-root dressing.

The cross-block orthogonality relations ($\delta_{j,j'}$) in
Eq.~\ref{eq:cpsv16_source_joint_isometry} hold for direct sums of distinct
normal blocks with idempotent transfer maps.\footnote{Formal theorem:
\leanid{MPSTensor.exists_residualIsometryFamily_of_isTransferIdempotent_directSum}.}
In CPSV canonical-form data, every block carries a unit-modulus weight and is
itself a renormalization fixed point. Selecting one representative per phase
class yields simultaneous trace-normalized square-root forms satisfying all
diagonal and off-diagonal residual isometry equations.\footnote{Formal theorem:
\leanid{MPSTensor.CPSVCanonicalFormData.BNTRefinement.exists_residualIsometryFamily_of_isTransferIdempotent},
wrapped by
\leanid{MPSTensor.IsCPSVCanonicalForm.exists_bntRefinement_residualIsometryFamily_of_isTransferIdempotent}.}

The repeated-copy clause of Theorem~\cpsvref{thm:charact-MPS} in
Ref.~\cite{Cirac2016MPDO_arXiv} remains underspecified: the displayed isometry
omits $q$, and the text refers to an unconstructed physical direct sum. The
tensor in Eq.~\ref{eq:cpsv16_literal_virtual_block_example} shows that a
literal virtual block diagonal violates both Eq.~(\cpsvref{AA=A}) and
transfer-map idempotence. The formal development records the exact criterion
for literal blocks, while the independent-phase assertion remains open
pending a clarified source statement.

\bibliographystyle{alpha}
\bibliography{references}

\end{document}
